%=====================================================================
\chapter{Structural Equation Modelling and Latent Constructs}\label{ch:sem}
%=====================================================================
\locator{4}{Part~4 --- Measurement, Data and Analysis}

\begin{outcomes}
By the end of this chapter you will be able to:
\begin{tight}
  \item \textbf{Explain} what a latent variable is and why measurement error
    matters for estimating relationships.
  \item \textbf{Separate} the measurement model from the structural model.
  \item \textbf{Choose} between covariance-based and partial least squares SEM
    on defensible grounds.
  \item \textbf{Assess} convergent and discriminant validity, including HTMT.
  \item \textbf{Detect} and address common method bias.
  \item \textbf{Report} an SEM study to the standard reviewers expect.
\end{tight}
\end{outcomes}

\begin{prereq}
\chref{measurement} for validity and reliability --- SEM is measurement theory
made operational. \chref{survey} if your data comes from a questionnaire, which
it usually will. \chref{regression} for the modelling ideas SEM extends.
\end{prereq}

\begin{keyterms}
latent variable $\bullet$ indicator $\bullet$ measurement model $\bullet$
structural model $\bullet$ reflective and formative indicators $\bullet$
confirmatory factor analysis $\bullet$ CB-SEM $\bullet$ PLS-SEM $\bullet$
composite reliability $\bullet$ average variance extracted $\bullet$
Fornell--Larcker $\bullet$ HTMT $\bullet$ common method bias $\bullet$ fit index
\end{keyterms}

\section{Why This Chapter Exists for an IT Course}

A large share of MS IT theses study technology acceptance, user satisfaction,
system quality or adoption intention. These are latent constructs measured by
questionnaire, related to one another in a model --- which is precisely what
structural equation modelling is for. Done well it is powerful; done as it
commonly is, it produces impressive diagrams supporting conclusions the data do
not carry.

\begin{definitionbox}[Latent variable]
A \term{latent variable} is a construct that cannot be observed directly and is
inferred from multiple observed \term{indicators}.

\medskip
``Perceived usefulness'' is latent; the four questionnaire items measuring it
are indicators. SEM estimates relationships \emph{between the latent
constructs}, explicitly modelling the fact that each indicator measures its
construct imperfectly.
\end{definitionbox}

\begin{conceptbox}[What SEM buys over regression on scale means]
Averaging four items and running a regression treats the average as if it were
the construct, measured without error. It is not, and the consequence is
systematic: measurement error \emph{attenuates} estimated relationships, biasing
them towards zero.

\medskip
SEM models the error explicitly and estimates the relationship between the
error-free constructs. For a scale with reliability 0.7, the attenuation is
substantial --- so this is not a refinement, it is a correction.

\medskip
The cost is that SEM requires more data, more assumptions, and considerably
more reporting. Do not use it because it looks sophisticated; use it because
your constructs are latent and measured with error.
\end{conceptbox}

\section{Two Models in One}

\begin{center}
\small
\begin{tabular}{@{}L{3.6cm}L{10.7cm}@{}}
\toprule
\textbf{Measurement model} & How each latent construct relates to its
                             indicators. Assessed \emph{first}, and the whole
                             analysis stops here if it fails \\
\textbf{Structural model}  & How the latent constructs relate to each other.
                             The hypotheses live here \\
\bottomrule
\end{tabular}
\end{center}

\begin{alertbox}[Assess the measurement model first, and be willing to stop]
If the indicators do not measure the constructs adequately, the structural
paths between them are meaningless --- you would be estimating relationships
between things you have not measured.

\medskip
The order is non-negotiable, and so is the willingness to act on it. A thesis
that reports poor discriminant validity and then proceeds to interpret path
coefficients anyway has told the examiner exactly what to ask about.
\end{alertbox}

\subsection{Reflective and formative indicators}

\begin{center}
\small
\begin{tabular}{@{}L{2.4cm}L{5.4cm}L{6.5cm}@{}}
\toprule
 & \textbf{Reflective} & \textbf{Formative} \\
\midrule
Direction  & Construct $\rightarrow$ indicators
           & Indicators $\rightarrow$ construct \\
Meaning    & Indicators are effects of the construct; they should correlate
             highly and are interchangeable
           & Indicators are components that constitute the construct; they need
             not correlate \\
Example    & Satisfaction: four items all reflecting how satisfied someone is
           & Socioeconomic status: income, education, occupation. Raising
             income raises status; it does not raise education \\
Assessment & Reliability, AVE, discriminant validity
           & Collinearity among indicators, and the significance of each
             indicator's weight \\
\bottomrule
\end{tabular}
\end{center}

Misspecifying the direction is a common and serious error: applying reflective
criteria to a formative construct will make a perfectly good construct look
invalid, and dropping ``poorly loading'' formative indicators removes part of
the construct's definition.

\section{CB-SEM or PLS-SEM?}

\begin{center}
\small
\begin{longtable}{@{}L{2.8cm}L{5.3cm}L{5.7cm}@{}}
\toprule
 & \textbf{CB-SEM} (\texttt{lavaan}) & \textbf{PLS-SEM} (\texttt{seminr}) \\
\midrule
\endfirsthead
\toprule
 & \textbf{CB-SEM} & \textbf{PLS-SEM} \\
\midrule
\endhead
\bottomrule
\endfoot
Estimates
  & Parameters reproducing the observed covariance matrix
  & Composite scores maximising explained variance in the outcomes \\
Purpose
  & Theory \emph{testing} and confirmation
  & Prediction and theory \emph{development} \\
Sample size
  & Larger; a common rule of thumb is 10 observations per estimated parameter
  & Smaller samples tolerated --- though the ``10 times the largest number of
    paths'' rule is widely criticised as far too permissive \\
Distribution
  & Assumes multivariate normality for maximum likelihood
  & Distribution-free; inference by bootstrapping \\
Global fit
  & Yes: $\chi^2$, CFI, TLI, RMSEA, SRMR
  & No true global fit test. SRMR is reported as an approximate check \\
Formative constructs
  & Awkward; requires additional identification constraints
  & Handles them naturally \\
\end{longtable}
\end{center}

\begin{pitfall}[PLS-SEM is not a way to escape a small sample]
It is chosen for that reason constantly, and the justification does not hold.
PLS tolerates smaller samples than CB-SEM, but small-sample PLS estimates are
still imprecise, and the bootstrap intervals will show it if you report them.

\medskip
Choose PLS because your model contains formative constructs, because your aim
is prediction, or because the theory is still being developed. Choose CB-SEM
because you are testing an established theory and want a global fit test. Then
say which reason applied, in one sentence, in the methods chapter. ``PLS-SEM
was used due to the small sample size'' is the sentence reviewers have learned
to distrust.
\end{pitfall}

\section{Assessing the Measurement Model}

\begin{center}
\small
\begin{longtable}{@{}L{3.6cm}L{3.6cm}L{6.6cm}@{}}
\toprule
\textbf{Criterion} & \textbf{Conventional threshold} & \textbf{What it checks} \\
\midrule
\endfirsthead
\toprule
\textbf{Criterion} & \textbf{Threshold} & \textbf{What it checks} \\
\midrule
\endhead
\bottomrule
\endfoot
Indicator loading
  & $\ge 0.70$
  & Each indicator relates strongly to its construct \\
Composite reliability
  & 0.70 to 0.95
  & Internal consistency. \emph{Above} 0.95 suggests redundant items, not an
    excellent scale \\
Average variance extracted (AVE)
  & $\ge 0.50$
  & The construct explains more than half the variance in its indicators \\
Fornell--Larcker
  & $\sqrt{\text{AVE}}$ exceeds correlations with other constructs
  & Discriminant validity. Now known to have poor sensitivity \\
HTMT
  & $< 0.85$, or $< 0.90$ for conceptually similar constructs
  & Discriminant validity, and substantially better at detecting its
    absence~\cite{henseler2015new} \\
\end{longtable}
\end{center}

\begin{conceptbox}[Report HTMT, and understand what failing it means]
Fornell--Larcker was the standard for two decades and has been shown to miss
discriminant validity problems in exactly the situations where they matter.
HTMT --- the heterotrait-monotrait ratio of correlations --- detects them
reliably, and reviewers in information systems now expect it.

\medskip
Failing HTMT means two of your constructs are not empirically distinguishable.
Your respondents did not separate them, whatever the theory says. The honest
responses are to merge the constructs and rewrite the model, or to redesign the
items so they diverge --- not to report the failure and interpret the path
between them anyway.
\end{conceptbox}

\section{Common Method Bias}

When one respondent supplies both the predictors and the outcome, in one
questionnaire, at one sitting, some of the observed correlation may be produced
by the method rather than by the constructs~\cite{podsakoff2003common}. Since
this describes nearly every acceptance study ever run, it deserves attention
rather than a formulaic paragraph.

\begin{center}
\small
\begin{tabular}{@{}L{3.4cm}L{10.9cm}@{}}
\toprule
\textbf{Procedural remedies} (better) & Separate the measurement of predictor
                                        and outcome in time or source;
                                        guarantee anonymity; counterbalance
                                        item order; avoid ambiguous items;
                                        use different response formats for
                                        different constructs \\
\textbf{Statistical detection} (weaker) & Harman's single-factor test --- widely
                                        used, widely criticised as
                                        insensitive; a marker variable
                                        approach; a common latent factor in the
                                        model \\
\bottomrule
\end{tabular}
\end{center}

\begin{alertbox}[Harman's single-factor test proves almost nothing]
It is reported in a large fraction of IS theses, usually as a single sentence
concluding that common method bias is not a concern. The test has very low
power to detect bias, so passing it is weak evidence, and the sentence is
performing reassurance rather than analysis.

\medskip
Design the bias out procedurally --- temporal separation between predictor and
outcome measurement is the single most effective step --- and then report what
you did in the design, not only what you tested afterwards. That is a stronger
answer, and it is available to you only before data collection.
\end{alertbox}

\begin{lstlisting}[language=Rlang, caption={A PLS-SEM assessment in the order that
matters: measurement model first, structural model only if it passes.}]
library(seminr)

measurements <- constructs(
  reflective("PU",  multi_items("PU",  1:4)),   # perceived usefulness
  reflective("PEOU",multi_items("PEOU",1:4)),   # perceived ease of use
  reflective("INT", multi_items("INT", 1:3))    # behavioural intention
)

structure <- relationships(
  paths(from = c("PU", "PEOU"), to = "INT"),
  paths(from = "PEOU",          to = "PU")
)

m <- estimate_pls(data = survey, measurements, structure)
s <- summary(m)

# 1. Measurement model. Stop here if it fails.
s$loadings; s$reliability          # loadings, rhoC, AVE
s$validity$htmt                    # discriminant validity -- the one to report

# 2. Structural model, with bootstrap intervals rather than p-values alone.
b <- bootstrap_model(m, nboot = 5000)
summary(b)$bootstrapped_paths
s$paths                            # R-squared per endogenous construct
\end{lstlisting}

\section{Reporting}

\begin{center}
\small
\begin{tabular}{@{}L{4.0cm}L{10.3cm}@{}}
\toprule
Why SEM, and which variant & The substantive reason, not the sample size \\
Every item                 & In an appendix, with its source or its adaptation
                             stated (\chref{survey}) \\
Measurement model          & Loadings, composite reliability, AVE, HTMT matrix
                             --- as tables, in full \\
Any items dropped          & Which, and why. Dropping items to reach a
                             threshold is a data-driven decision and must be
                             disclosed \\
Structural model           & Path coefficients with bootstrap confidence
                             intervals, $R^2$, and effect sizes \\
Fit (CB-SEM)               & $\chi^2$ with df, CFI, TLI, RMSEA, SRMR --- all of
                             them, not the flattering subset \\
Common method bias         & What was done procedurally, then what was tested \\
Data and code              & Deposited (\chref{reproducibility}) \\
\bottomrule
\end{tabular}
\end{center}

\begin{pitfall}[Four things examiners ask about SEM theses]
\begin{tight}
  \item \textbf{``Why these paths?''} Every path must come from theory stated
    in advance. A model with paths between everything is not a hypothesis.
  \item \textbf{``What did you drop?''} Items removed to improve AVE, and
    modification indices followed to improve fit, are data-driven decisions.
    Fit obtained that way is not evidence for the theory.
  \item \textbf{``Is this causal?''} Cross-sectional SEM estimates
    associations. Arrows in the diagram are the theory's claims, not the
    design's warrant (\chref{quasiexp}).
  \item \textbf{``Would a different model fit as well?''} Usually yes.
    Equivalent models with reversed paths often fit identically, and
    acknowledging this is far better than being shown it.
\end{tight}
\end{pitfall}

\begin{traditions}{latent variable modelling}
\tradEMP{Its natural extension: measurement error is modelled rather than
ignored, and theory is tested against a covariance structure.}
\tradDSR{Used to evaluate an artefact's perceived qualities and their relation
to intention to use --- formative naturalistic evaluation, and weak evidence of
actual use.}
\tradQUAL{Would object that constructs like ``perceived usefulness'' are
imposed on respondents rather than derived from them, and that an exploratory
sequential design (\chref{mixed}) should establish the constructs first.}
\tradFORM{Identification --- whether parameters can be uniquely determined ---
is a formal question with a formal answer, and an unidentified model is not a
statistical problem but an algebraic one.}
\end{traditions}

\begin{lab}[Assess a measurement model]
Using an open dataset, or the pilot data from your own instrument
(\chref{survey}):

\begin{tightnum}
  \item Specify three constructs with at least three reflective indicators
    each, and the paths between them, from theory.
  \item Estimate the measurement model. Report loadings, composite reliability
    and AVE.
  \item Compute both Fornell--Larcker and HTMT. Do they agree? If not, which do
    you trust and why?
  \item Only if the measurement model passes, estimate the structural model
    with 5,000 bootstrap resamples and report intervals.
  \item Write the paragraph you would put in the thesis about common method
    bias --- describing design steps first.
\end{tightnum}
\end{lab}

\begin{checkpoint}
\begin{tightnum}
  \item Why does measurement error bias regression coefficients towards zero,
    and how does SEM address it?
  \item Give a formative construct from IT research and explain why reflective
    assessment criteria would misjudge it.
  \item Composite reliability of 0.97 is reported as excellent. What is the
    concern?
  \item Two constructs fail HTMT. What are your two legitimate options?
\end{tightnum}
\end{checkpoint}

\begin{chaptersummary}
\begin{tight}
  \item Latent constructs measured with error attenuate relationships; SEM
    corrects for this by modelling the error.
  \item Measurement model first, structural model only if it passes.
  \item Reflective and formative indicators require different assessment;
    misspecifying the direction invalidates the assessment.
  \item Choose CB-SEM or PLS-SEM for a substantive reason. Small sample size is
    not one.
  \item Report HTMT, not only Fornell--Larcker. Failing it means merging
    constructs or redesigning items.
  \item Design common method bias out procedurally; Harman's single-factor test
    is weak evidence.
  \item Disclose every dropped item and every modification index followed. Fit
    obtained by data-driven revision is not evidence for the theory.
\end{tight}
\end{chaptersummary}

\begin{reviewq}
  \item Technology acceptance models are among the most replicated in IS and
    among the most criticised as theoretically thin. Assess that criticism, and
    say what a genuinely novel contribution using UTAUT2 would have to add.
  \item Equivalent models can fit identically while asserting opposite causal
    directions. What does this imply about the evidential value of a
    well-fitting cross-sectional SEM, and what design would break the tie?
  \item PLS-SEM's sample size rules of thumb have been shown to be far too
    permissive. Why do they persist, and what should replace them?
  \item \reg{PhD Regs 2024}{\S14.2(v)} requires that ``quantitative research
    proposals must include a valid statistical design for data analysis''. What
    would a valid design section for an SEM study contain that most contain
    now?
\end{reviewq}
